%% ══════════════════════════════════════════════════════════
%%  RAA — Collapsing Flat Waves (Robert Rehnig)
%%  Sienas teksts: "Calm Down But Keep Breathing"  (latviešu)
%% ══════════════════════════════════════════════════════════

\documentclass[a4paper,10.5pt]{article}

\usepackage[top=2.2cm,bottom=2cm,left=2.8cm,right=2.8cm]{geometry}
\usepackage{xcolor}
\usepackage[T1]{fontenc}
\usepackage[utf8]{inputenc}
\usepackage{graphicx}
\usepackage{parskip}
\usepackage{microtype}
\usepackage{lmodern}
\usepackage{enumitem}
\usepackage{hyperref}

\hypersetup{hidelinks}

\definecolor{raared}{HTML}{8B0000}
\definecolor{raadark}{HTML}{0D0D0D}
\definecolor{mygray}{HTML}{555555}

\pagestyle{empty}
\setlength{\parskip}{6pt}

\newcommand{\ruledtitle}[1]{%
  \vspace{4pt}%
  {\color{raared}\bfseries #1}\\[-4pt]%
  {\color{raared}\rule{\linewidth}{0.6pt}}%
  \vspace{2pt}%
}

\begin{document}

%% ── Galvene ───────────────────────────────────────────────
\noindent
\includegraphics[height=1.4cm]{img/raa_logo_v2.png}

\vspace{0.55cm}

%% ── Izstādes nosaukums ────────────────────────────────────
{\small\color{mygray}\textsc{Izstāde}}\\[1pt]
{\LARGE\bfseries Collapsing Flat Waves}\\[1pt]
{\normalsize\color{mygray}Robert Rehnig \ \textbullet\ 24.09.--14.10.2026 \ \textbullet\ RAA.SPACE, Rīga}

\vspace{0.55cm}

%% ── Darba nosaukums ───────────────────────────────────────
{\small\color{mygray}\textsc{Skaņu instalācija}}\\[1pt]
{\Large\bfseries\color{raared}``Calm Down But Keep Breathing''}

\vspace{0.3cm}

%% ── Darba apraksts ────────────────────────────────────────
Šī reaģējošā telpiskā skaņu instalācija atveido akustiski elpojošu galerijas telpu -- vai izdomātu, abstraktu būtni. Apgaismotajā telpas centrā staru formēšanas skaļruņu masīvs veido it kā kolonnu, kurā dzirdama ieelpotā gaisa skaņa ķermeņa iekšienē, bet ap duci pie sienām piestiprinātu skaļruņu projicē izelpotā gaisa skaņu, kas izplatās iedomātā, lielākā ēkā. Izmērītais skaļums galerijas iekšienē un ārpusē ietekmē šo akustisko būtni, un tā reaģē, uz laiku apklustot.

\vspace{0.15cm}
{\color{raared}\rule{\linewidth}{0.8pt}}
\vspace{0.35cm}

%% ── Par mākslinieku ───────────────────────────────────────
\ruledtitle{Par Mākslinieku}

Robert Rehnig (dz. 1977) ir Leipcigā dzīvojošs skaņu mākslinieks un komponists. Viņš strādā audiogrāmatu, skatuves mūzikas un skaņu instalāciju mākslas jomās. Viņa sērija ``Souvenirs entomologiques'' saņēmusi Vācijas audiogrāmatu balvu. Viņa darbi prezentēti (izlase): Singuhr-H\"orgalerie Berlīnē, Festspielhaus Hellerau Drēzdenē, Kunstverein Hamburg, ZixP festivālā Leipcigā, Mid-Autumn Harvest Moon Festival Monreālā, Maison du G\'eant Salins-les-Bains, Patarei Tallinā, festivālā Punto de Encuentro Valensijā, sound:frame festivālā K\"unstlerhaus Vīnē un Fr\"uhjahrstage f\"ur zeitgen\"ossische Musik Veimārā. Viņš pasniedz SeaM (Elektroakustiskās mūzikas un skaņu mākslas studijā, Veimāras Mūzikas augstskolā ``Franz Liszt'') un Leipcigas Mūzikas un teātra augstskolas ``Felix Mendelssohn Bartholdy'' elektroniskās mūzikas studijā. Pašlaik viņš ir doktorants Veimāras Bauhaus universitātē, kur pēta klusuma psihofiziku un fenomenoloģiju.

\vfill

%% ── Kājene: RAA rekvizīti ─────────────────────────────────
{\color{raared}\rule{\linewidth}{0.8pt}}
\vspace{0.3cm}

\noindent
\begin{minipage}[t]{0.72\textwidth}
{\small
\textbf{RAA.SPACE}\\[3pt]
Instagram \ \href{https://www.instagram.com/raa_riga/}{@raa\_riga}\\
Mājaslapa \ \ \href{https://raa.space}{raa.space}\\
Izstādes lapa \ \href{https://raa.space/events/collapsing-flat-waves}{raa.space/events/collapsing-flat-waves}\\
Instalācijas ieraksts \ \href{https://www.instagram.com/p/DdRzAJ_I8d0/}{instagram.com/p/DdRzAJ\_I8d0}\\[6pt]
\color{mygray}Matīsa iela 8, Rīga, LV-1001, Latvija
}
\end{minipage}%
\begin{minipage}[t]{0.28\textwidth}
\raggedleft
\includegraphics[width=2.6cm]{img/qr-collapsing-flat-waves.pdf}
\end{minipage}

\end{document}
